\section{Recovering the leading geometry from intrinsic marks}
\label{sec:calibration}
All arclength variables in this section refer to one representative of
\eqref{eq:intrinsicmeasure}. The quantities used for calibration are
unchanged when that representative is reversed.

\begin{proposition}[Onset and quadratic calibration]\label{prop:calibration}
The local onset is $T_0=2g$. The conditional rescaled endpoint law
converges to the probability measure on $\R^2$ with density
\begin{equation}\label{eq:ellipselimit}
 \frac{a_b}{\pi}(1-\tfrac12Y^tH_bY)_+\dd Y.
\end{equation}
In particular, its covariance is $H_b^{-1}/3$ and
\begin{equation}\label{eq:betavalue}
 \beta_b=\frac{2g}{3c_b},\qquad
 \beta_b^+:=\lim_{d\downarrow0}
 \frac{\int(s_-+s_+)^2\dd\mu_{b,2g+d}}{dP_b(d)}
 =\frac{2g}{3(c_b-c_{1-b}^{-1})}.
\end{equation}
Consequently \eqref{eq:calibrationmain} holds. The plus moment is an
optional compatibility check; it is not needed by the inverse.
\end{proposition}
\begin{proof}
The reduced action has a strict minimum $2g$ at the normal orbit by
Lemma~\ref{lem:flux}. On a sufficiently small fixed patch no competing
minimum occurs. There is no event for $T\le2g$, whereas the positive
twist and the open residual interval give positive mass for every
sufficiently small $T>2g$. This proves the onset statement.

Let $y=\sqrt d\,Y$ in the physical integral. Its active set stays in a
bounded ellipse, because $E_b(y)\ge c|y|^2$ near the minimum. Taylor's
formula gives $E_b(\sqrt dY)/d\to E_{b,0}(Y)$,
$b_b(\sqrt dY)\to1$ and
$\mathfrak a_b(\sqrt dY_i)/\sqrt d\to Y_i$ uniformly there. The constant $a_b$ in
\eqref{eq:ellipselimit} is $\sqrt{\det H_b}$, as in \eqref{eq:H}.
After division by $P_b(d)$, dominated convergence gives
\eqref{eq:ellipselimit}, including its bounded quadratic moments.

Under $X=H_b^{1/2}Y$, the normalizing integral is
$\int_{|X|<\sqrt2}(1-|X|^2/2)\dd X=\pi$.
Its second moment in either coordinate is $\pi/3$, and its mixed
moment is zero. Thus the covariance is $H_b^{-1}/3$.
The vector $(1,-1)^t$ is an eigenvector of $H_b$ with eigenvalue
$c_b/g$, and $(1,1)^t$ has eigenvalue $(c_b-c_{1-b}^{-1})/g$.
This proves \eqref{eq:betavalue} without knowing either curvature in
advance. Solving its first formula gives
$\kappa_b=2/(3\beta_b)-1/g$. Finally
$\sinh(2\gamma)=2\sqrt{c_0c_1(c_0c_1-1)}$ and the leading physical
probability coefficient give the area formula.
\end{proof}

Only intrinsic boundary distances enter this computation. The formula is
not an inversion of the support of an embedded endpoint cloud: no
longitudinal boundary coordinate is observed. The two curvatures are
separated by the two itineraries even at equal curvature.

\begin{corollary}[Leading-coordinate Jacobian]\label{cor:leadingchart}
On $g,\kappa_0,\kappa_1,A>0$, the map
$(g,\kappa_0,\kappa_1,A)\mapsto(T_0,\beta_0,\beta_1,\alpha)$
has an analytic inverse. In this order its determinant is
\begin{equation}\label{eq:leadingdet}
 -\frac{8\alpha g^4}{9A c_0^2c_1^2}\ne0.
\end{equation}
Together with the higher-jet blocks it gives an analytic $2K$-dimensional
finite-jet coordinate map on a chosen orientation cover.
\end{corollary}
\begin{proof}
After the first row, eliminate the $g$ derivatives in the two
$\beta$ rows. Their remaining diagonal derivatives are
$-2g^2/(3c_b^2)$. The last diagonal derivative is
$\partial_A\alpha=-\alpha/A$. Multiplication with
$\partial_gT_0=2$ gives \eqref{eq:leadingdet}. The number of higher
coordinates is $2(K-2)$.
\end{proof}
